“The continuous photometric and distance-layer realisation of this Clock, including magnification bias μlens, isotropic weight Wiso, and the integrated χ(z), dL(z), lives in the companion repository:
https://github.com/TheAstrographer/JCR-Kerr-Gordon_Lensing-Metric_Starter-Pack”

Retaining the Dual Architecture  
Mathematical-Numerical Precision and Lensing-Photometric Grounding

The Strictly Mathematical / Numerical Layer  

This layer is defined by exact discrete relations and high-precision continuous realisations:

Tick lattice: \(\varepsilon = 10^{-9}\), \(N_{\rm today} = 10^{9}\)
\(\chi_{\rm micro} = \varepsilon N\), \(a(N) = \exp(\chi_{\rm micro})\), \(z(N) = 1/a(N) - 1\)
Angular bridge \(\psi = 0.1503378808\), phase slip \(\Delta\phi_{\rm torque} = 1.72113420759\)
Exact identity \(\sin(\Delta\phi_{\rm torque}) = \cos\psi\)
Torque amplitude locked at \(3.170\) km s\(^{-1}\) Mpc\(^{-1}\)
Machine-precision \(H_{\rm eff}(0) = 73.170000\)
Exact mapping of every tick \(N\) onto the model redshift \(z(N)\)
Rapidly decaying suppression that leaves high-\(z\) expansion nearly untouched

This layer lives in pure numerical precision. Every identity closes to machine accuracy. It supplies the kinematic skeleton and the local \(H_0\) target.

2. The Lensing-Photometric Layer  
(Kerr–Gordon Lensing Metric Starter Pack)

This layer is grounded in the observational and geometric consequences of light propagation:

Isotropic conformal weight \(W_{\rm iso}(\beta,z)\) with \(\beta = 0.05\)
Explicit magnification bias  
  \[
  \mu_{\rm lens}(z) = 1 + 0.05\,\frac{z}{1+z}
  \]
Comoving and luminosity distances obtained by integrating \(c/H_{\rm eff}(z)\)
Sachs equation modified by the vortex / torque stress-energy contribution \(\pi_{\mu\nu}\)
Photometric bias that directly affects measured magnitudes and therefore the distance ladder and supernova Hubble diagram

This layer is not an abstract numerical identity; it is the physical interface between the geometric torque and actual astronomical measurements—magnification, photometric shifts, and the way flux is recorded. It is the part that can be compared with real photometric catalogues and lensing statistics.

The mathematical-numerical layer guarantees internal coherence and the precise local value \(73.17\).
The lensing-photometric layer supplies the observable consequences.

One layer is the clockwork (exact ticks, phases, and torque amplitude). The other is the face of the clock (how light is magnified and photometrically registered).

The rapidly decaying torque belongs to both: it is calculated with numerical precision in the mathematical layer and then inserted into the distance and magnification integrals of the photometric layer. The mapping \(N\to z(N)\) likewise remains exact in the mathematical layer while providing the redshift coordinate at which photometric quantities are evaluated.
